% ============================================
% LUC-20 — O Arquivo Aberto (edição revista R1075)
% ============================================
\fragtitle{luc}{L}{LUC}{20}{O Arquivo Aberto}
\markboth{LUC-20 --- O Arquivo Aberto}{}

\noindent \textit{Arquivo Público Mineiro --- 14 de maio de 2026, 2h10}

\noindent O arquivo à noite era outro lugar. As luzes de emergência deixavam os corredores num amarelo baixo, e as estantes altas pareciam paredes de um labirinto. Lúcia passou o cartão na porta da sala de documentos restritos. A fechadura piscou verde. Ela entrou.

\noindent Ronaldo estava lá.

\noindent Sentado na cadeira dela, com o diário aberto sobre a mesa. Não se assustou quando ela entrou. Levantou os olhos devagar, como se já soubesse a hora.

\noindent ``Você não devia ter vindo'', ele disse.

\noindent ``Você não devia estar aqui.''

\noindent ``Eu vim fechar a sala.''

\noindent ``A sala está aberta.''

\noindent Ele olhou para o diário. ``Eu quis ler. Só uma página. Para saber se era verdade.''

\noindent ``E era?''

\noindent Ronaldo não respondeu. Havia algo diferente no rosto dele. A pele parecia mais fina, os ossos mais visíveis. O olho direito estava amarelo. Não o amarelo de reflexo. O amarelo de dentro.

\noindent Lúcia se aproximou da mesa. O diário estava aberto na página 62. A página estava coberta de uma letra apertada, miúda, que ela conhecia: era a letra de Joaquim. Mas havia outra letra nas margens, em tinta mais escura. Uma letra de médico.

\noindent ``O que você escreveu?'' ela perguntou.

\noindent ``Nada.''

\noindent ``Ronaldo.''

\noindent ``Eu juro.'' Ele levantou as mãos. ``Eu só li. A página estava assim quando eu abri.''

\noindent Lúcia olhou para as margens. Havia uma frase repetida três vezes, em três caligrafias diferentes: ``A fila não diminuiu.'' A primeira era de Joaquim. A segunda era de Oliveira. A terceira era dela.

\noindent Ela sentiu o chão se mover. O chão ficou firme. O corpo se moveu. O coração bateu uma vez, forte, e depois pareceu parar. O ar do arquivo ficou denso, com cheiro de terra e papel molhado.

\noindent ``Você escreveu isso'', ela disse.

\noindent ``Não.''

\noindent ``Você escreveu isso agora.'' Ela apontou para a terceira linha. ``Esta letra é a minha. Eu não escrevi. Mas é a minha.''

\noindent Ronaldo olhou para a página. Depois para ela. Havia medo nos olhos dele, mas não o medo de quem descobre. O medo de quem confirma.

\noindent ``Eu não fui o primeiro'', ele disse. ``Meu pai leu. Depois eu li. Depois você. A fila não pede escolha. Alguém sempre segura a vez. Alguém sempre segura a vez.''

\noindent ``E você segurou a minha.''

\noindent ``Eu segurei o seu nome.'' Ele engoliu. ``Quando o diário voltou, em 2019, eu podia ter queimado. Podia ter enterrado. Podia ter deixado na caixa. Mas eu abri. E a primeira frase que eu li foi o meu nome. Não o do Joaquim. O meu. Escrito com a minha letra. Eu tinha nove anos. Eu não sabia escrever assim.''

\noindent Lúcia se lembrou da fotografia da mãe dele, do armário queimado, do caderno preto. Havia uma cadeia ali que não começava em 1853. Começava antes. Começava em cada pessoa que decidia passar adiante.

\noindent ``Por que você não queimou?'' ela perguntou.

\noindent ``Porque eu tinha medo de que, se queimasse, ele viesse atrás de mim. E porque...'' Ele parou. ``Porque uma parte de mim queria saber o que tinha no fim.''

\noindent Lúcia olhou para o diário. A página 62 parecia iluminada por dentro, uma luz amarela fraca, como uma lâmpada velha. Ela estendeu a mão. Ronaldo segurou seu pulso.

\noindent ``Não'', ele disse. ``Se você ler, não volta.''

\noindent ``Eu já li.''

\noindent ``Não a última página.''

\noindent ``Não existe última página.''

\noindent Ele abriu a boca para responder. As luzes do arquivo piscaram. Uma vez. Duas. A sala ficou escura. No escuro, Lúcia ouviu a respiração de Ronaldo, rápida, curta, e depois outra respiração, mais lenta, vindo da estante. Havia alguém ali. Ou alguma coisa.

\noindent ``Ronaldo'', ela disse.

\noindent Ninguém respondeu.

\noindent A luz de emergência voltou. Ronaldo estava no chão, encostado na estante, com a mão no peito. Os olhos abertos. O olho direito completamente amarelo. Ele não se mexia.

\noindent Lúcia se ajoelhou. Tocou o rosto dele. Estava frio. No bolso do paletó, havia um envelope. Ela tirou. Dentro, uma folha com uma frase datilografada:

\noindent ``Não fui eu que escrevi o diário. Fui eu que decidi quem ia ler.''

\noindent Ela olhou para Ronaldo. Depois para o diário. Depois para a página 62. A luz amarela se apagou. No lugar, havia uma linha nova, escrita com a letra dela:

\noindent ``Agora é a minha vez.''

\noindent Lúcia fechou o diário. Guardou no pano preto. Colocou na caixa de metal. Fechou o cadeado. Saiu do arquivo sem olhar para trás. Na recepção, o segurança dormia. Ela deixou a caixa sobre o balcão, com uma etiqueta:

\noindent ``Devolver ao paciente 1.263, quando aparecer.''

\noindent O segurança acordou quando ela abriu a porta. Perguntou se estava tudo bem.

\noindent ``Está'', ela disse. ``Só estou devolvendo o que não é meu.''

\noindent Ele olhou para a caixa. Depois para ela. Depois para a porta.

\noindent ``A senhora esqueceu a chave.''

\noindent Lúcia olhou para a mão. A chave estava lá. Quente.

\noindent ``Não'', ela disse. ``A chave não é minha.''

\noindent Saiu. A porta fechou atrás dela. No arquivo, a caixa de metal ficou sobre o balcão, com o diário dentro, esperando. O segurança voltou a dormir. Sonhou com uma vala. Quando acordou, o cheiro doce estava no ar.

\noindent\textit{\textbf{↪ Rotas:} \rota{LUC-19} \rotasep \rota{CONT-01} \rotasep \rota{MEM-27}}
